\section{Implementierung}
\label{sec:implementierung}

\subsection{Containerarchitektur und Persistenz}

Die fünf Container kapseln inkompatible CUDA- und Python-Abhängigkeiten. SAM
läuft in einem modernen CUDA-12.6-Container, COLMAP in einem eigenen SfM-
Container, STS in CUDA~12.1, SuGaR in CUDA~11.8 und DGtal/GDAL CPU-seitig.
Bind-Mounts\footnote{Bind-Mounts sind Verzeichniseinbindungen, bei denen ein
Host-Ordner direkt an eine Stelle im Container gemountet wird; Dateien werden
dabei nicht in die Container-Ebene kopiert, sondern beide Seiten greifen auf
dieselben Dateisystemobjekte zu. Dadurch bleiben Artefakte nach Beendigung
des Containers auf dem Host erhalten.} schreiben alle relevanten Artefakte in
den gemeinsamen Host-Arbeitsbereich.

\begin{table}[H]
\centering
\small
\begin{tabularx}{\textwidth}{p{1.4cm}p{3.4cm}XX}
\toprule
Stufe & Kernsoftware & Eingabe & Ausgabe\\
\midrule
A & SAM~3.1 & Video, Textprompt & Frames und hierarchische Masken\\
B & COLMAP & unmaskierte Frames & Kameras, Sparse-Modell, ideale Szene\\
C & STS & ideale Szene und Masken & objektannotierte 3D-Gaussians\\
D & SuGaR-Fork/Open3D & Objekt-Gaussians und Kameras & Coarse-Mesh,
optional Refined-Mesh\\
E & DGtal/GDAL & kompatibles OBJ & lokale Centerline, B-Spline, CSV/GeoJSON\\
\bottomrule
\end{tabularx}
\caption{Verantwortlichkeiten und zentrale Datenübergaben der fünf Container.}
\label{tab:containerrollen}
\end{table}

Die Master-Skripte übernehmen die numerische Host-UID und -GID, das ist die
Benutzer- und Gruppenkennung des ausführenden Nutzers auf dem Host. Dadurch
gehören neu erzeugte Dateien auf dem Host dem eigenen Benutzer und nicht dem
Container-Root; bind-gemountete Dateien bleiben so auch auf Hochschul- und
Domänensystemen lesbar und bearbeitbar.
Run-Tags und Ausgabeordner werden vor unbeabsichtigtem Überschreiben geschützt.
Manifeste, Parameter und Logs bleiben außerhalb kurzlebiger Container erhalten.

\subsection{Mask-aware SuGaR-Fork und Codeänderungen}

Der für die Vergleichsroute verwendete SuGaR-Stand ist kein unverändertes
Vanilla-SuGaR. Er liegt als projektinterner Fork unter
\path{third_party/SuGaR} und wurde als Git-Submodul in die Pipeline eingebunden.
Die Erweiterungen sind funktionale Codeänderungen, nicht nur andere
Startparameter.

Der Grund dafür liegt in einem grundlegenden Zielkonflikt der
Standardimplementierung: In frühen Entwicklungstests mit unmodifiziertem
Vanilla-SuGaR führte der unveränderte Full-Image-RGB-Loss dazu, dass das
System eine bereits auf das Rohr reduzierte Punktwolke weiterhin gegen den
gesamten Boden- und Wandhintergrund optimierte. Wie in
Abbildung~\ref{fig:pa-panel-mesh} (Teilbild~2) dokumentiert, führte dieser
Widerspruch zwischen 3D-Geometrie (nur Rohr) und 2D-Supervision (gesamtes
Foto) zu massiven Floatern und seitlich in den Raum ragenden Dreiecksflügeln.

Um 3D-Eingabe und 2D-Bildsupervision auf dieselbe Objektdefinition zu
vereinheitlichen, wurde der SuGaR-Quellcode um einen maskierten
Optimierungsverlust ($\mathcal{L}_{\mathrm{RGB}}^{M}$,
Gleichung~\ref{eq:masked-rgb-loss}) erweitert. Für die standardisierte
Versuchsmatrix war dieser maskenbewusste Fork von Beginn an als feste
Vergleichsbasis integriert. Die in früheren Matrixläufen beobachteten Ausfälle
dieser Route waren folglich keine Schwäche des maskierten Trainings, sondern
ein rein technischer Fehler im nachgelagerten Renderpfad (Import- und
Tensorlayoutfehler beim Erzeugen der Testansichten), der vor der finalen
12er-Folgematrix behoben wurde.

\begin{table}[H]
\centering
\small
\begin{tabularx}{\textwidth}{p{2.4cm}p{4.1cm}X}
\toprule
Forkstand & Geänderte Bereiche & Funktionale Änderung\\
\midrule
\texttt{48bbfdd} & \path{sugar_utils/mask_utils.py},
\path{sugar_trainers/}, \path{sugar_extractors/}, \path{train.py} &
Einführung eines \texttt{SemanticMaskProvider}; maskierter RGB-/L1-/DSSIM-
Loss in den Coarse-Trainern und im Refinement; maskierte
Depth-Normal-Consistency; konfigurierbare Maskenstufen und Dilatationen;
maskenkonsistentes UV-Baking; zusätzliche Mesh- und Trainingsargumente.\\
\texttt{8f4f7a2} & \path{sugar_extractors/coarse_mesh.py}, \path{train.py} &
Schutz vor zu kleinen Punktwolken und ungültigen Normalen bei der
Poisson-Rekonstruktion; Korrektur der Coarse-Zählerprüfung für geladene STS-
Checkpoints oberhalb von 6999.\\
\texttt{e254000} & \path{extract_mesh.py},
\path{sugar_extractors/coarse_mesh.py} &
Diagnoseoption für Tiefen aus dem Gaussian-Rasterizer gegenüber der
projizierten Diamond-Mesh-Tiefe.\\
\texttt{ecda7ef} & \path{extract_mesh.py},
\path{sugar_extractors/coarse_mesh.py}, \path{sugar_scene/sugar_model.py} &
Konfigurierbare Surface-Sampling-Parameter und explizite Behandlung der
Background-Mesh-Komponente.\\
\texttt{eca4ea1} / Arbeitsstand & \path{gaussian_splatting/scene/dataset_readers.py},
\path{sugar_scene/gs_model.py} &
Normalisierte feste Eval-Splits und Schutz gegen nicht gematchte Kameras; im
aktuellen Arbeitsbaum wird die Split-Kardinalität zusätzlich explizit geprüft.\\
\bottomrule
\end{tabularx}
\caption{Funktionale Änderungen des maskenbewussten SuGaR-Forks.}
\label{tab:sugar-fork-aenderungen}
\end{table}

Der aktuelle Parent-Gitlink zeigt auf \texttt{ecda7ef}; der ausgecheckte
Entwicklungszweig steht auf \texttt{eca4ea1}. Die zusätzlichen Änderungen in
\path{dataset_readers.py} und \path{gs_model.py} sind zum Zeitpunkt dieser
Fassung noch lokale Arbeitsbaumänderungen. Vor der Abgabe müssen Parent-
Gitlink, Fork-Commit und Docker-Parameter \texttt{SUGAR\_REF} auf einen
identischen, reproduzierbaren Stand gebracht werden. Die vollständige
Commit-Historie und der lokale Diff gehören in den digitalen Anlagenindex.

Für die Projektarbeit ist diese Forkentwicklung wissenschaftlich relevant,
weil dadurch zwei Wirkorte derselben SAM-Semantik getrennt geprüft werden:
STS selektiert die Objekt-Gaussians, die Coarse-Route verwendet die Maske für
ihre Bildsupervision und UV nutzt sie nur für die spätere Farbakkumulation.
Eine UV-Maske kann daher das Coarse-Mesh nicht verändern; eine RGB-Maske kann
hingegen die trainierte Gaussian-Geometrie beeinflussen.

\subsection{Mehrfache Bild- und Maskenablagen}

Die scheinbar häufigen Duplikate entstehen überwiegend durch unterschiedliche
Datenverträge, nicht dadurch, dass SAM dieselbe Inferenz mehrfach als
notwendigen Funktionsschritt ausführt. Die aktuelle Ablage lässt sich wie folgt
trennen:

\begin{table}[H]
\centering
\small
\begin{tabularx}{\textwidth}{p{4.0cm}p{2.3cm}X}
\toprule
Ablage & Physische Kopie? & Zweck und Bewertung\\
\midrule
\path{data/02_frames/} & Ja & Von SAM erzeugte Arbeitsframes. Sie bilden den
gemeinsamen Rohbildsatz für SAM und COLMAP und sind funktional notwendig.\\
\path{data/04_sfm/undistorted/images/} & Ja & Durch
\texttt{image\_undistorter} neu abgebildete ideale Bilder. Diese Kopie ist im
aktuellen Domänenkonzept notwendig, weil STS nicht auf den verzerrten Rohpixeln,
sondern auf der idealen PINHOLE-Szene arbeitet.\\
\path{data/05_3dgs/images} & Nein, Symlink & STS benötigt den stabilen
\texttt{source\_path/images}-Vertrag. Der Link verweist auf die ideale Szene
und erzeugt keinen zweiten Bildbestand.\\
\path{data/03_masks/_attempts/} & Ja & Prompt-Fallbacks werden getrennt
gespeichert, damit bei OOM, leerer Maske oder ungeeignetem Prompt der beste
Versuch ausgewählt werden kann. Nach der Auswahl sind diese Dateien für die
Funktion nicht mehr nötig, aber als Diagnose wertvoll.\\
\path{data/03_masks/} und \path{data/05_3dgs/multiview\_masks\_*} & Ja &
Die flachen Masken in \path{data/03_masks/} sind die Primärdaten. Die
\path{multiview_masks_*}-Ordner unter \path{05_3dgs} sind STS-spezifische
Kompatibilitätskopien (Zuordnung zu Objekt-ID 000); STS liest daraus. Die
\path{*_merged}-Verzeichnisse entstehen intern während der STS-Initialisierung
und sind flüchtige Arbeitsdaten. Die Hierarchie (\texttt{default},
\texttt{middle}, \texttt{small}) ist semantisch nicht drei Inferenzläufe,
sondern eine unveränderte Maske plus morphologische Ableitungen (Dilatation
und Erosion).\\
\path{data/10_runs/} & Ja & Matrixläufe werden vollständig archiviert und vor
dem nächsten Versuch vom Live-Workspace getrennt. Das ist Reproduzierbarkeit,
keine zusätzliche Verarbeitung im Einzelrun.\\
\bottomrule
\end{tabularx}
\caption{Funktionale Gründe für mehrere Bild- und Maskenablagen.}
\label{tab:bildablagen}
\end{table}

Die wichtigste physische Bildkopie entsteht somit nicht durch eine unnötige
SAM-Optimierung, sondern durch die geometrische Neuberechnung in die ideale
Bilddomäne. Sie ließe sich später speichereffizienter organisieren, etwa durch
einen Cache oder eine kontrollierte gemeinsame Workspace-Struktur, in der die
idealen Bilder ohne zweite physische Kopie bereitgestellt werden. Eine solche
Optimierung ist für die Funktion des aktuellen
Workflows nicht elementar: Würde sie die idealen Bilder oder den stabilen
STS-Pfadvertrag entfernen, wären Kameras, Masken und Pixelkoordinaten nicht
mehr konsistent.

Speicher- und I/O-Optimierungen sind dennoch möglich. Die Promptversuche in
\path{_attempts} können nach erfolgreicher Auswahl archiviert und aus dem
Arbeitsbereich entfernt werden. Die STS-Kompatibilitätskopien können langfristig
regenerierbar statt dauerhaft gespeichert werden. Der Symlink unter
\path{data/05_3dgs/images} sollte dagegen erhalten bleiben. Diese Maßnahmen
reduzieren Speicherbedarf und Kopierzeit, verändern aber nicht die zentrale
Rekonstruktionslogik und sind deshalb nachrangig gegenüber korrekter
Masken-/Kameradomäne, COLMAP-Registrierung und validierten Quality-Gates.

\subsection{CLI, Autopilot und Replay}

Die Pipeline verwendet einen geführten Vollpfad sowie einen Replay-Pfad.
\path{run_pipeline.sh} mit \texttt{--from <step>} startet diesen Replay. Er
ist kein Fortsetzen innerhalb eines abgebrochenen Trainings, sondern ein
Neustart an einer
Schrittgrenze mit bereits vorhandenen, zueinander passenden Artefakten.

Die Konzentration auf ein zentrales CLI-Skript ist eine bewusste
Scopeentscheidung. Zunächst soll eine robuste, automatisierbare Anwendung
existieren; eine zusätzliche Bedienoberfläche würde eigene Zustands-, Upload-
und Prozessfragen eröffnen. Der Autopilot verwendet unter anderem 5 FPS,
$1280\times720$ Pixel, CRF 18, \texttt{OPENCV}, 4096 SIFT-Merkmale,
Overlap 15, deaktiviertes Guided Matching, 7000 STS-Iterationen und Route A.
Das Kameramodell \texttt{OPENCV} wurde auf Grundlage der Matrixauswertung in
Abschnitt~\ref{sec:produktionskonfiguration} von der historischen
\texttt{SIMPLE\_RADIAL}-Baseline zum Standard erhoben. Die Parameter werden
protokolliert und können im interaktiven Modus erklärt beziehungsweise
überschrieben werden.

Der interaktive Vollauf fragt nach der Autopilot-Entscheidung zusätzlich nach
einem Lauf-Preset, das die Auflösung des Arbeitsvideos und damit die
Frame-/STS-/SuGaR-Auflösung festlegt: 720p ($1280\times720$, Standard),
qHD ($960\times540$) oder low ($640\times360$). Bei qHD und low werden die
Zielauflösungsfragen im Video-Preprocessing fest vorbesetzt; bei 720p bleiben
alle Eingabefragen erhalten. Die Abfrage weist die aus den Matrixläufen
gemittelten Laufzeiten der Produktionskonfiguration (Route A, 5 FPS,
\texttt{OPENCV}) aus: rund 34 Minuten für 720p, 27 Minuten für qHD und
19 Minuten für low. Diese Werte umfassen die archivierten Schrittzeiten von
STS bis Postprocessing; SAM~3.1, COLMAP, Rendering und Metrikberechnung sind
nicht vollständig enthalten. Im Autopilot-Modus wird automatisch das
720p-Preset mit allen Standardeingaben gewählt; die gewählte Auflösung wird
im Run-Report protokolliert.

\subsection{Quality-Gates und Datenverträge}

Qualitätsrelevante Implementierungsentscheidungen sind:

\begin{itemize}
  \item COLMAP verwendet unmaskierte Bilder, damit Hintergrund- und
        Ankermerkmale für SfM und GCP erhalten bleiben.
  \item Matrixrunner und Inline-Vollauf prüfen die Maskenabdeckung vor STS;
        der Guard bricht bei mindestens 30\,\% leeren oder fehlenden
        \texttt{middle}-Masken ab.
  \item Matrixrunner und Inline-Vollauf warpen die Rohmasken nach COLMAP in
        dieselbe ideale Bilddomäne wie die STS-Bilder. Die Rohmasken bleiben
        unter \path{data/03_masks/} erhalten; die gewarpten Masken liegen unter
        \path{data/03_masks_ideal/} und werden an STS sowie SuGaR übergeben.
        Der Matrixrunner erzeugt zusätzlich einen festen Eval-Split aus den
        nichtleeren idealen Masken. Ein Replay ab \texttt{sts} verwendet bereits
        bereitgestellte ideale Masken; ein Replay ab \texttt{colmap} besitzt
        diese automatische Warp-Stufe weiterhin nicht und benötigt daher einen
        passenden idealen Maskensatz.
  \item Die SuGaR-Kameranamen werden als normalisierte Stems verglichen; leere
        Kamera- oder Eval-Listen werden mit klaren Fehlern abgelehnt.
  \item Die Bildauswertung ist ausschließlich objektmaskiert. Full-frame-Werte
        würden bei einem object-only Splat überwiegend den nicht rekonstruierten
        Hintergrund bewerten.
\end{itemize}

Zusätzlich prüfen Replay- und Matrixpfad, ob ein vollständiges COLMAP-Modell
vorliegt, die Objektfilterung nicht alle Gaussians entfernt, die Poisson-
Eingabe gültige Punkte und Normalen enthält und die für einen Replay
notwendigen Artefakte aus demselben Frame-, Kamera- und Iterationsstand stammen.
Der Matrixrunner archiviert jeden Lauf, bevor der Live-Arbeitsbereich für die
nächste Konfiguration bereinigt wird.

\subsection{STS-Objektfilterung und geometrischer Eingang}

Der vollständige STS-Checkpoint bleibt unverändert erhalten. Ein separates
Skript selektiert die konfigurierte Objekt-ID und kann Primitive unterhalb der
Opazität 0,01 beziehungsweise sehr dunkle Primitive unterhalb eines
Farbgrenzwerts von 0,08 entfernen. Fällt die Filterung leer aus, verhindert ein
Quality-Gate die unbemerkte Übergabe. Die verbliebenen Datensätze werden in
einer getrennten Kopie auf $\alpha=0{,}999999$ gesetzt. Dadurch bleibt die
Originalwolke für Diagnose und Vergleich verfügbar.

\subsection{Historische Zwischenstände der STS- und Objektfilterungsroute}

Die folgenden Entwicklungsstände dokumentieren, welche Einstellungen und
Fehlerbefunde schließlich zu den Parametern der aktuellen Route geführt haben.
Sie sind historische Belege aus unterschiedlichen Entwicklungsphasen und
bewusst nicht als quantitative Ergebnisreihe zu lesen.

\begin{figure}[H]
  \centering
  \includegraphics[width=\textwidth]{pa_panel_sts_entwicklung.png}
  \caption{Historische STS-Entwicklung: von ersten Fehlversuchen mit
  invertierter Auswahl des segmentierten Trainings (Fehlorientierung im
  Rendering) über Floaters und fehlerhafte Zwischenstände bis zu
  verlässlichen Objekt-Splats. Die Bildausschnitte stammen aus verschiedenen
  Entwicklungsständen (Juni bis Juli 2026) und untermauern die schrittweise
  Fixierung der STS-/Rendering-Datenverträge.}
  \label{fig:pa-panel-sts}
\end{figure}

\begin{figure}[H]
  \centering
  \includegraphics[width=\textwidth]{pa_panel_mesh_filterung.png}
  \caption{Historische Meshgewinnung und Objektfilterung: vom ersten
  herausgeschnittenen Objekt über schlecht segmentierte Zwischenstände und
  Floaters bis zum segmentierten Refined-Mesh der aktuellen Route.}
  \label{fig:pa-panel-mesh}
\end{figure}

Das STS-Panel (Abbildung~\ref{fig:pa-panel-sts}) zeigt den historischen
Entwicklungsweg des Renderings: Die Teilbilder 2 bis 4 dokumentieren Floater
und die Probleme ohne die später eingeführte Farbfilterung; Teilbild 5 belegt,
dass auch bei passender Segmentierung eine schlechte Aufnahme noch immer
problematische Splats erzeugen kann. Das Mesh-Panel
(Abbildung~\ref{fig:pa-panel-mesh}) dokumentiert entsprechend die Meshroute:
Teilbild 1 und 2 zeigen, wie sich SuGaR vor dem Fork verhielt und warum der
maskenbasierte Verlust nötig wurde. Teilbild 3 entstand aus einem fälschlich
übergebenen Splat, der direkt in ein Mesh umgewandelt wurde und extrem viele
Floater hinterließ. Teilbild 4 stammt von der zuerst erprobten
SuGaR-Coarse-Variante mit Zielzähler 15.000 und bewusst deaktivierter
DN-Consistency -- die lang gezogenen Strukturen entstanden dort jedoch durch
die falsche Übergabe und die fehlende Maskenbasiertheit beziehungsweise die
RGB-lastige Supervision, nicht durch die Route selbst. Die Teilbilder 5 und 6
zeigen einen Vergleich perfekt segmentierter Splats, Teilbild 7 die
zugehörige Punktwolke. Insgesamt belegen beide Panels, dass ein
funktionierendes STS-Rendering nicht selbstverständlich ist und welche
Fehlerbilder die heutigen Quality-Gates ausschließen.
\subsection{Splat-Rendering und Generierung der Bildmetriken}

Die Bildmetriken entstehen nicht im Nachgang aus dem extrahierten Dreiecksnetz,
sondern werden über automatisierte Rendering-Skripte direkt aus den
3D-Gaussian-Checkpoints auf dem festen Test-Split (\path{eval_frames.txt})
erzeugt. Die Pipeline unterscheidet dabei drei getrennte Metrik-Dateien:

\begin{enumerate}
  \item \textbf{\path{metrics/sts_masked.json} (Original-STS / Route A):}\\
        Nach Abschluss des STS-Curriculums (7000 Iterationen) wird der Standard-3DGS-
        Checkpoint \path{point_cloud/iteration_7000/point_cloud.ply} geladen. Das
        Skript \path{render_standard_splat} rendert die Ansichten des festen
        Eval-Splits -- also Kameras, deren Bilder nicht im Training enthalten
        waren; anschließend berechnet \path{evaluate_masked_splat_metrics.py} die
        objektmaskierten Werte für PSNR, SSIM und LPIPS gegen diese
        Ground-Truth-Bilder.
        Da Route~A das Mesh direkt aus diesem STS-Splat ohne weiteres Training
        extrahiert, bildet diese Datei die photometrische Qualitätsmetrik für Route~A.
  \item \textbf{\path{metrics/sugar_coarse_masked.json} (SuGaR-Coarse-Vergleich):}\\
        In der SuGaR-Vergleichsroute optimiert das Coarse-Training ($c=9000$) die
        Gaussians mit maskiertem RGB-Loss und Normalenkonsistenz und speichert das
        Ergebnis als PyTorch-Checkpoint unter
        \path{data/sugar_output/<run_tag>/coarse/9000.pt}. Das Skript
        \path{src/python/render_sugar_checkpoint.py} lädt dieses \texttt{state\_dict}
        in das \texttt{SuGaR}-Modell in Container~D und rendert denselben
        Test-Kamerasplit. Die anschließende maskierte Auswertung schreibt
        \path{sugar_coarse_masked.json}. Sie misst die Ansichtsqualität der durch
        SuGaR veränderten Gaussian-Geometrie.
  \item \textbf{\path{metrics/sugar_refined_masked.json} (SuGaR-Refined):}\\
        Diese Datei ist für das Rendering eines verfeinerten, mesh-gebundenen Splats
        (\path{refined.ply}) vorgesehen. Wurde im jeweiligen Lauf kein
        \path{refined.ply} exportiert, schreibt der Matrixrunner ein explizites
        \texttt{\{"status":"skipped","reason":"no refined.ply exported"\}}.
\end{enumerate}

Für das korrekte SuGaR-Rendering waren zwei technische Korrekturen in
\path{render_sugar_checkpoint.py} notwendig: Erstens muss der SuGaR-Quellordner
\path{/opt/sugar} vor dem Import von \texttt{sugar\_scene} explizit in den
Python-Suchpfad (\texttt{sys.path}) eingefügt werden. Zweitens liefert der
SuGaR-Rasterizer ein HWC-Tensorlayout, während \texttt{torchvision.utils.save\_image}
ein CHW-Layout erwartet; die Hilfsfunktion \texttt{as\_chw()} normalisiert die
Kanaldimension, um Tensorformfehler zu verhindern.

\subsection{Meshrouten und kompatibler Export}

Die Produktionsroute A führt keine SuGaR-Coarse-Optimierung durch. Sie nutzt
Surface-Sampling, Poisson-Rekonstruktion, Dichtebereinigung, Decimation und
Vertexprojektion. Die Legacy-Route \texttt{sugar\_coarse} wird für eine
kontrollierte Gegenüberstellung separat ausgeführt.

Route A speichert das Coarse-PLY und konvertiert es für die bestehende
Centerline-Schnittstelle zusätzlich in \texttt{refined.obj}. Dieser
Kompatibilitätsname bedeutet ausdrücklich nicht, dass ein SuGaR-Refinement
ausgeführt wurde. Die Vergleichsroute schreibt Coarse-, Refined- und
gegebenenfalls Texturartefakte getrennt.

\subsection{Centerline und Export}

Container E führt den \texttt{single}-Centerline-Modus mit einer Voxelgröße von
0,1 und einer Mindestpfadlänge von 0,75 aus. Die Eckensegmentierung ist
deaktiviert; die lokale Kurve wird als geklemmter uniformer B-Spline Grad 10
ausgegeben. Fehlt eine Transformationsmatrix, werden globale Ausgaben deutlich
als Translation-Fallback bezeichnet. Die detaillierten GeoJSON- und
Postprocessdateien befinden sich im digitalen Anlagenteil.

\subsection{Ausgewählte Robustheitsmaßnahmen}

Die Entwicklung enthielt mehrere negative Ergebnisse. Für den Haupttext sind
nicht die jeweiligen Stacktraces, sondern die daraus abgeleiteten
Schnittstellenregeln relevant. Tabelle~\ref{tab:robustheit-kern} fasst die
wichtigsten Beispiele zusammen; die ausführlichere Chronologie befindet sich
in Tabelle~\ref{tab:fehlerchronik} im Anhang.

\begin{table}[H]
\centering
\small
\begin{tabularx}{\textwidth}{p{2.5cm}XXX}
\toprule
Problem & Ursache & Korrektur & Bedeutung\\
\midrule
Leerer SuGaR-Eval-Split & Dateiendung und Kamerastem unterschieden sich &
Normalisierung und Kardinalitätsprüfung & Ein vorhandener Split ist erst nach
erfolgreichem Matching gültig.\\
Unvollständige Matrix & Unterprozesse konsumierten TSV-stdin &
stdin-unabhängige Schleifen & Geplante und tatsächlich ausgeführte Läufe müssen
getrennt geprüft werden.\\
Leere Eval-Maske & Globaler Coverage-Guard war weniger streng als der feste
Split & Split erst nach dem Warp aus nichtleeren Masken & Eval-Eignung ist ein
eigener Datenvertrag.\\
SuGaR-Renderfehler & Importpfad und HWC/CHW-Layout & expliziter Pfad und
Layoutkonvertierung & STS-Zwischenmetriken dürfen nicht als SuGaR-Ergebnis
gedeutet werden.\\
Verästeltes Skelett & Medialflächen und kurze Stachel-Äste & Durchmesserpfad im
\texttt{single}-Modus & Netzwerkobjekte bleiben außerhalb des Scopes.\\
\bottomrule
\end{tabularx}
\caption{Kernprobleme und daraus abgeleitete Robustheitsmaßnahmen.}
\label{tab:robustheit-kern}
\end{table}
